% =============================================================
% BODY (ES) – GENERAL TEMPLATE CONTENT
% =============================================================

\section{Overview}
\vspace{2mm}

Este documento utiliza una plantilla bilingüe (ES/EN) orientada a documentación técnica y personal. El objetivo es mantener una estructura modular, consistente y fácil de mantener en el tiempo.

\vspace{2mm}


\subsection{Qué Se Edita}
\vspace{2mm}

\begin{itemize}
  \item \textbf{Metadatos:} editar el bloque \texttt{METADATA} (título, categoría, versión, estado, fecha, autor).
  \item \textbf{Contenido:} escribir en \texttt{lang/es/body.tex}.
  \item \textbf{Referencias:} completar \texttt{src/references.tex} (bibliografía manual al final).
  \item \textbf{Assets:} guardar recursos en \texttt{assets/figures/}, \texttt{assets/diagrams/} y \texttt{assets/tables/}.
\end{itemize}

\vspace{2mm}


\section{Contenido Principal}
\vspace{2mm}

Se recomienda mantener párrafos breves, y separar por subsecciones cuando el tema crezca (una idea principal por subsección).

\vspace{2mm}


\subsection{Notas, Tips Y Advertencias}
\vspace{2mm}

\begin{noteBox}
\boxlabel{Nota}
Bloque para decisiones, contexto relevante o aclaraciones.
\end{noteBox}

\vspace{2mm}

\begin{tipBox}
\boxlabel{Tip}
Bloque para recomendaciones prácticas, buenas prácticas o shortcuts seguros.
\end{tipBox}

\vspace{2mm}

\begin{warnBox}
\boxlabel{Advertencia}
Bloque para riesgos, restricciones o pasos propensos a errores.
\end{warnBox}

\vspace{2mm}


\subsection{Checklist}
\vspace{2mm}

\begin{itemize}
  \item Punto a verificar.
  \item Supuesto o decisión tomada.
  \item Resultado esperado.
  \item Próximo paso.
\end{itemize}

\vspace{2mm}


\section{Ejemplos}
\vspace{2mm}

\subsection{Código}
\vspace{2mm}

\begin{lstlisting}
# Example
print("Hello, world!")
\end{lstlisting}

\vspace{2mm}


\subsection{Figura (Snippet)}
\vspace{2mm}

\begin{lstlisting}
\begin{figure}[htbp]
  \centering
  \includegraphics[width=0.80\linewidth]{assets/figures/your-figure.png}
  \caption{Descripcion breve de la figura.}
\end{figure}
\end{lstlisting}

\vspace{2mm}


\subsection{Diagrama / Compuertas (TikZ)}
\vspace{2mm}

\begin{center}
\begin{tikzpicture}[circuit logic US, scale=1.0, transform shape]
  \node[and gate, draw, logic gate inputs=nn] (G) at (0,0) {};
  \draw (-1.2,0.2) -- (G.input 1);
  \draw (-1.2,-0.2) -- (G.input 2);
  \draw (G.output) -- (1.2,0);
\end{tikzpicture}
\end{center}

\vspace{2mm}